\chapter{Hecke Theory for \texorpdfstring{$\GL(2)$}{GL(2)}}

\label{ch:6}

Let $F$ denote a global field and $\A$ its ring of adeles. If $\psi$ is a grossencharacter of $F$ let $\bL^{2}(G_F \backslash G_{\A}, \psi)$ denote the Hilbert space of measurable functions $\phi$ on $\GL(2, F) \backslash \GL(2, \A)$ satisfying

\begin{equation}\label{eq:6.1}
\phi(zg) = \psi(z)\phi(g) \quad \text{for all } z \in Z_{\A}\tag{6.1}
\end{equation}

and

\begin{equation}\label{eq:6.2}
\int_{Z_{\A}G_F \backslash G_{\A}} |\phi(g)|^2 \dd g < \infty.\tag{6.2}
\end{equation}

As in Section~\ref{ch:5}, let $\bR^{\psi}(g)$ denote the unitary representation of $\bG_{\A}$ given in $\bL^2(G_F \backslash G_{\A}, \psi)$ by right translation.

\begin{definition}\label{def:6.1}

An irreducible unitary representation $\pi$ of $G_{\A}$ will be called an automorphic form for $\GL(2)$ if it occurs in some $R^{\psi}$. More precisely, suppose

\begin{equation}\label{eq:6.3}
R^{\psi}(g) = \left(\bigoplus_j \pi^j\right) \oplus \int \pi^s \dd s.\tag{6.3}
\end{equation}

Then $\pi$ is an automorphic form if it is equivalent to some $\pi^j$ or $\pi^s$. It is a cusp form if it is equivalent to some $\pi^j$ (and not 1-dimensional).

\end{definition}

Clearly certain cusp forms in the above sense correspond to classical cusp forms in $S_k(N,\psi)$ and these forms have associated to them Dirichlet series which by Hecke and Weil essentially characterize them. (Cf. Sections~\ref{ch:1}.C and~\ref{ch:1}.D.)

It might be expected therefore that all generalized cusp forms have associated to them certain Dirichlet series and that these Dirichlet series characterize them in the following sense. Suppose an ``L-function'' can be attached to a given irreducible unitary representation $\pi$ of $G_{\A}$ and suppose this function satisfies certain prescribed analytic conditions. Then $\pi$ must be a cusp form. That this is indeed the case is the main assertion of Jacquet-Langlands' theory.

The purpose of this section is to describe Jacquet-Langlands' characterization of cusp forms and to explain some implications of their theory for the classical theory of forms.

\section{Hecke Theory for \texorpdfstring{$\GL(1)$}{GL(1)}}\label{sec:6a}

This subsection and the next are included primarily to help motivate the group theoretic techniques of Jacquet and Langlands.

In this subsection only, $G_{\A}$ will denote the idele group $\A^{\times} = \GL(1,\A)$ of a number field $F$. By analogy with Definition~\ref{def:6.1}, we shall understand an automorphic form on $\GL(1)$ to be any irreducible unitary representation of $G_{\A} (= \GL(1,\A) = \A^{\times})$ which occurs in the decomposition of the natural representation $R(g)$ of $G_{\A}$ in $\bL^2(G_F \backslash G_{\A}) (= \bL^2(F^{\times} \backslash \A^{\times}))$. But

\begin{equation}\label{eq:6.4}
R(g) = \oplus \int_{G_F \backslash G_{\A}} \psi(g) \dd \psi.\tag{6.4}
\end{equation}

Therefore automorphic forms on $\GL(1)$ are simply grossencharacters of $F$, i.e.\ unitary characters of $\A^{\times} = G_{\A}$ trivial on $G_F = F^{\times}$, and Hecke's theory of L-functions with grossencharacter is simply a theory of Dirichlet series attached to automorphic forms for $\GL(1)$!

This theory may be described as follows. Suppose $\psi$ is any grossencharacter of $F$. Then there are characters $\psi_v$ of $(F_v)^{\times}$ (for each place $v$ of $F$) such that

\begin{equation}\label{eq:6.5}
\psi = \prod_v \psi_v\tag{6.5}
\end{equation}

with $\psi_v$ unramified for almost every $v$.

Let $B$ denote the (finite) set of places of $F$ ``ramified'' for $\psi$, i.e.\ such that $\psi_v$ is ramified. For each $v$ not in $B$, put

\begin{equation}\label{eq:6.6}
\chi(\mathfrak{v}) = \psi_v(\tilde{\omega}_v)\tag{6.6}
\end{equation}

where $\tilde{\omega}_v$ denotes a local uniformizing parameter at $v$ (a generator of the prime ideal of $O_v$). Note that $\chi(\mathfrak{v})$ does not depend on the choice of $\tilde{\omega}$. Indeed $\psi_v$ is trivial on the units of $F_v$. Therefore $\chi$ is well defined by multiplicativity on the group of $F$-ideals relatively prime to $B$ (the ``bad'' primes) and the character $\chi$ is the Hecke character associated to the grossencharacter $\psi$.

The Hecke L-series associated to $\chi$ (or $\psi$) is defined by the Dirichlet series

\begin{equation}\label{eq:6.7}
L(s, \chi) = \sum \frac{\chi(\mathfrak{A})}{N(\mathfrak{A})^s},\tag{6.7}
\end{equation}

where the summation is taken over all $F$-ideals relatively prime to $B$, or, equivalently, by the Euler product

\begin{equation}\label{eq:6.8}
L(s, \chi) = \prod_{v \notin B} \left(1 - \frac{\chi(\mathfrak{v})}{(N(\mathfrak{v}))^s}\right)^{-1} = \prod_{v \notin B} \left(1 - \frac{\psi_v(\omega_v)}{|\omega_v|^{-s}}\right)^{-1}.\tag{6.8}
\end{equation}

Although it is easy to check that this series (or product) converges absolutely uniformly to the right of the line $\operatorname{Re}(s) = 1$ it is not so easy to prove Hecke's principal result concerning this series.

\begin{theorem}\label{thm:6.2}

\begin{enumerate}[(i)]

\item $L(s,\chi)$ has an analytic continuation as a meromorphic function of $s$. If $\chi$ is not ``trivial'' then $L(s,\chi)$ is actually entire; otherwise $L(s,\chi)$ has a simple pole at $s = 1$ as its only singularity;

\item $L(s,\chi)$ satisfies a functional equation of the following type: there is a constant $A$, and a gamma factor $\Gamma(s,\chi)$ (depending only on the infinite primes of $F$) such that if

\[
R(s,\chi) = s(s-1)A^s\Gamma(s,\chi)L(s,\chi)
\]

then $R(s,\chi)$ is entire and

\begin{equation}\label{eq:6.9}
R(1 - s, \chi^{-1}) = W(\chi) R(s, \chi)\tag{6.9}
\end{equation}

with $W(\chi)$ a constant of modulus 1.

\end{enumerate}

\end{theorem}

Hecke's proof of this result uses the classical theory of theta functions and represents a great achievement.

The arithmetic significance of Theorem~\ref{thm:6.2} is well-known. Indeed if $\chi$ is trivial, $L(s,\chi)$ is the Dedekind zeta-function of $F$. Thus, taking $F = \Q$, Theorem 2 gives (in particular) the analytic continuation and functional equation of the Riemann zeta function.

A drawback of Hecke's proof is that it does not lay bare the group theoretic interpretation of the result. This was left to be done by Tate in his justly celebrated thesis. Since Tate's method has since proved to be germinal to many developments in number theory (let alone Jacquet-Langlands theory) it seems worthwhile to sketch it below.

Let $S(\A)$ denote the Schwartz-Bruhat functions on $\A$. These functions are linear combinations of functions of the form $\prod_v f_v$ where

\begin{enumerate}[(i)]

\item $f_v$ is in the usual Schwartz space of $\R$ or $\C$ if $v$ is archimedean;

\item $f_v$ belongs to $S(F_v)$, the space of locally constant compactly supported functions on $F_v$ if $v$ is finite; and

\item $f_v$ is the characteristic function of the ring of integers of $F_v$ for almost every finite $v$.

\end{enumerate}

If $\tau$ is any fixed non-trivial character of $\A$ trivial on $F$, then the Fourier transform of any function $f$ in $S(\A)$ is given by

\begin{equation}\label{eq:6.10}
\hat{f}(\alpha) = \int_{\A} f(\beta) \tau(\alpha\beta) \dd\beta.\tag{6.10}
\end{equation}

Now suppose $\psi$ is an automorphic form on $G_{\A}$, i.e.\ $\psi$ is a grossencharacter of $F$. For each $f$ in $S(\A)$ the global zeta-function $\zeta(f, \psi, s)$ is defined by

\begin{equation}\label{eq:6.11}
\zeta(f, \psi, s) = \int_{G_{\A} = \A^{\times}} f(a)\psi(a) |a|^s \dd^{\times} a\tag{6.11}
\end{equation}

where $|a| = \prod_v |a_v|_v$ and $\dd^{\times} a = \prod_v \dd^{\times} a_v$. If $\psi = \prod_v \psi_v$ and $f = \prod_v f_v$ it is not difficult to show that for $\operatorname{Re}(s) > 1$,

\begin{equation}\label{eq:6.12}
\prod_v \int_{F_v} |f_v(a_v)| |a_v|_v^s |\dd^{\times} a_v| < \infty\tag{6.12}
\end{equation}

and

\begin{equation}\label{eq:6.13}
\zeta(f, \psi, s) = \prod_v \zeta(f_v, \psi_v, s),\tag{6.13}
\end{equation}

where $\zeta(f_v, \psi_v, s)$ denotes the local zeta-function

\begin{equation}\label{eq:6.14}
\zeta(f_v, \psi_v, s) = \int_{(F_v)^{\times} = G_v} f_v(a_v)\psi_v(a_v) |a_v|_v^s \dd^{\times} a_v.\tag{6.14}
\end{equation}

Therefore $\zeta(f, \psi, s)$ is easily seen to represent an analytic function of $s$ for $\operatorname{Re}(s) > 1$. What Tate proved was:

\begin{theorem}\label{thm:6.3}

The function $\zeta(f, \psi, s)$ has an analytic continuation to the whole $s$-plane, as an entire function if $\psi$ is non-trivial, and as a meromorphic function with simple poles at 0 and 1 otherwise. Furthermore, $\zeta(f, \psi, s)$ satisfies the functional equation

\begin{equation}\label{eq:6.15}
\zeta(f, \psi, s) = \zeta(\hat{f}, \psi^{-1}, 1 - s).\tag{6.15}
\end{equation}

\end{theorem}

\begin{corollary}\label{cor:6.4}

The function $\pi^{-s/2}\Gamma\left(\frac{s}{2}\right)\zeta(s)$ extends to a meromorphic function in $\C$ (with poles at 0 and 1) which is invariant under the change $s \to 1 - s$.

\end{corollary}

\begin{proof}

Take $F = \Q$, $\psi$ trivial, and $f = \prod_p f_p$, where $f_{\infty}(x) = e^{-\pi x^2}$, and $f_p$ is the characteristic function of $O_p$. Then (choosing $\tau$ as in Section~\ref{ch:3})

\[
\zeta(f_{\infty}, \psi_{\infty}, s) = \pi^{-s/2}\Gamma\left(\frac{s}{2}\right)
\]

and

\[
\zeta(f_p, \psi_p, s) = \int_{O_p} |\alpha|^s \dd^{\times}\alpha = (1 - p^{-s})^{-1}.
\]

Consequently

\[
\zeta(f, \psi, s) = \pi^{-s/2}\Gamma\left(\frac{s}{2}\right)\zeta(s)
\]

and the Corollary follows immediately from Theorem 6.4.

\end{proof}

\begin{corollary}\label{cor:6.5}

Theorem~\ref{thm:6.2} is valid.

\end{corollary}

\begin{proof}[Proof (Sketch)]

Using harmonic analysis on $G_v = (F_v)^{\times}$ it can be shown that the local zeta-functions themselves possess analytic continuations to the $\C$-plane and satisfy functional equations of the form

\begin{equation}\label{eq:6.16}
\zeta(f_v, \psi_v, s) = \rho(\psi_v, s)\zeta(\hat{f}_v, \psi_v^{-1}, 1 - s)\tag{6.16}
\end{equation}

where $\rho(\psi_v, s)$ is a meromorphic function of $s$ which does not depend on $f_v$. (The ``additive'' Fourier transform $\hat{f}_v$ is defined with respect to the character $\tau_v$ if $\tau = \Pi\tau_v$.)

A straightforward computation then shows that for properly chosen $f_v$,

\begin{equation}\label{eq:6.17}
\prod_{v \notin B_{\psi}} \zeta(f_v, \psi_v, s) = L(s, \chi) \prod_{v \notin B_{\psi}} \chi(\partial_v)^{-1} N(\partial_v)^{s - 1/2}\tag{6.17}
\end{equation}

where $B_{\psi}$ is the union of the set of archimedean primes and finite primes ramified for $\psi$ and $\partial_v$ is the different of $F_v$. But

\[
\zeta(f, \psi, s) = \prod_v \zeta(f_v, \psi_v, s).
\]

Therefore, by multiplying both sides of \eqref{eq:6.17} by $\prod_{v \in B_{\psi}} \zeta(f_v, \psi_v, s)$, it follows that $L(s, \chi)$ (like $\zeta(f, \psi, s)$ and $\prod_{v \in B_{\psi}} \zeta(f_v, \psi_v, s)$) possesses an analytic continuation to all of $\C$. Similarly the functional equation for $L(s, \chi)$ results from the functional equation for the global and local zeta-functions $\zeta(f, \psi, s)$. (The precise analytic behavior of $L(s, \chi)$ as well as an effective computation of $A$, $\Gamma(s, \chi)$ and $W(\chi)$ follows from a careful analysis of the factors $\rho(\psi_v, s)$.)

\end{proof}

\begin{remark}[\normalfont (Concerning the proof of Theorem~\ref{thm:6.3}.)]\label{rem:6.6}

A complete proof of Theorem~\ref{thm:6.3} is to be found in several places (see ``Notes and References''). We wish to emphasize that the proof is almost entirely Fourier analytic. More precisely, it represents a blend of (additive and multiplicative) harmonic analysis on the groups $M = M(1, F)$ and $G = \GL(1)$, the main ingredients being the Poisson summation formula, some generalized Mellin transforms, and a computation of the measure of the homogeneous space

\[
Z_{\infty}^{+} G_F \backslash G_{\A} = R_{+} F^{\times} \backslash \A^{\times}.
\]

We should also remark that nowhere in the proof of Theorem~\ref{thm:6.3} do the local zeta-functions $\zeta(f_v, \psi_v, s)$ play a real role. In fact these zeta-functions first significantly appear in the proof of Corollary~\ref{cor:6.5}. Thus in Tate's work the local zeta-functions are principally introduced to show that the theory is non-empty!

\end{remark}

\textbf{Concluding Remarks 6.7.} It seems to make sense to ask if there is a ``Converse'' to Hecke's theory for $\GL(1)$. In other words, suppose $\psi$ is a character of $G_{\A} = \A^{\times}$ such that $L(s, \chi)$ (defined as in 6.7) satisfies the analytic properties announced in Theorem~\ref{thm:6.2}. Is $\psi$ necessarily a grossencharacter, i.e.\ is $\chi$ necessarily trivial on $F^{\times}$?

\addtocounter{theorem}{1}

\section{Further Motivation}\label{sec:6b}

It is clearly indicated now that one should associate to every automorphic form $\pi$ an L-function with functional equation. Indeed this is what Hecke did for $\GL(1)$ and $\GL(2)$ in case $\pi$ corresponds to a classical holomorphic cusp form. It is also indicated that one should characterize the L-functions so obtained for general $\pi$ following Weil's Theorem~\ref{thm:1.12}. Both these goals are attained by Jacquet-Langlands' theory. They work in the generality of an arbitrary global field and their new methods involve the theory of group representations.

In Subsection C we shall sketch some principal results of Jacquet-Langlands' theory. Since their methods seem complex at first sight we shall attempt in the present subsection to provide motivation for them.

How does Hecke's theory attach a Dirichlet series to $f$ in $S_k(N, \psi)$? According to \eqref{eq:1.24}, one simply takes the Mellin transform of $f$ along the line $\{iy : y > 0\}$. That is,

\begin{equation}\label{eq:6.18}
L(s, f) = \int_0^{\infty} f(iy) y^{s-1} \dd y.\tag{6.18}
\end{equation}

(Strictly speaking, the Dirichlet series attached to $f$ is $D(s,f) = (2\pi)^s\Gamma(s)^{-1}L(s,f)$. However our interest is in the L-function $L(s,f)$, a simple modification of $D(s,f)$.)

To understand how one might generalize Hecke's construction it is helpful to rewrite \eqref{eq:6.18} in the framework of the adele group of $\GL(2)$.

As in Section~\ref{ch:3}, let

\begin{equation}\label{eq:6.19}
\phi_f(g) = f(g_{\infty}(i)) j(g_{\infty}, i)^{-k} \psi(k_0)\tag{6.19}
\end{equation}

denote the function on $G_{\A} = \GL(2, \A)$ corresponding to $f$ in $S_k(N, \psi)$. For simplicity, suppose that $\psi \equiv 1$ and $N = 1$. Then $\phi_f(g)$ is right invariant by $K_0 = \prod_{p < \infty} K_p$ (and, as always, left invariant by $G_{\Q}$). From \eqref{eq:6.19} it follows that (for $y > 0$)

\[
\phi_f\!\left(\begin{bmatrix} y & 0 \\ 0 & 1 \end{bmatrix}\right) = f(iy).
\]

Therefore \eqref{eq:6.18} is equivalent to the formula

\begin{equation}\label{eq:6.20}
L(s, f) = \int_{\Q^{\times} \backslash \A^{\times}} \phi\!\left(\begin{bmatrix} y & 0 \\ 0 & 1 \end{bmatrix}\right) |y|^s \dd^{\times} y\tag{6.20}
\end{equation}

(Recall that $\A^{\times} = \Q^{\times}\R_{+}(\prod_{\mathfrak{p} < \infty} O_{\mathfrak{p}}^{\times})$.)

To continue ``adelizing'' \eqref{eq:6.18} one must exploit the existence of a Fourier expansion for $f(z)$. According to Lemma~\ref{lem:3.6},

\begin{equation}\label{eq:6.21}
\begin{aligned}
\phi\!\left(\begin{bmatrix} y & 0 \\ 0 & 1 \end{bmatrix}\right) &= \sum_{\xi \in \Q^{\times}} \phi_{\xi}\!\left(\begin{bmatrix} y & 0 \\ 0 & 1 \end{bmatrix}\right)\\
&= \sum_{n=1}^{\infty} a_n e^{-2\pi ny} = f(iy)
\end{aligned}\tag{6.21}
\end{equation}

and

\begin{equation}\label{eq:6.22}
\phi\!\left(\begin{bmatrix} y & 0 \\ 0 & 1 \end{bmatrix}\right) = \sum_{\xi \in \Q^{\times}} W_{\phi}\!\left(\begin{bmatrix} \xi y & 0 \\ 0 & 1 \end{bmatrix}\right)\tag{6.22}
\end{equation}

where $W_{\phi}(g)$ is the K-finite $C^{\infty}$ function defined on $G_{\A}$ by

\begin{equation}\label{eq:6.23}
W_{\phi}(g) = \int \phi\!\left(\begin{bmatrix} 1 & x \\ 0 & 1 \end{bmatrix} g\right) \overline{\tau(x)} \dd x = \phi_1(g).\tag{6.23}
\end{equation}

Consequently, from \eqref{eq:6.20},

\[
L(s, f) = \int_{\Q^{\times} \backslash \A^{\times}} \sum_{\xi \in \Q^{\times}} W_{\phi}\!\left(\begin{bmatrix} y\xi & 0 \\ 0 & 1 \end{bmatrix}\right) |y|^s \dd y^{*},
\]

or

\begin{equation}\label{eq:6.24}
L(s, f) = \int_{\A^{\times}} W_{\phi}\!\left(\begin{bmatrix} y & 0 \\ 0 & 1 \end{bmatrix}\right) |y|^s \dd^{*} y.\tag{6.24}
\end{equation}

What \eqref{eq:6.24} says is that $L(s, f)$ is the adelic Mellin transform ``along $\begin{bmatrix} y & 0 \\ 0 & 1 \end{bmatrix}$,'' of the first Fourier coefficient of a special function $\phi_f$ in the space of $\pi_f$.

Suppose now that $\pi$ is an arbitrary constituent of some $R_0^{\psi}$. From the classical considerations above it is suggested that one might define $L(s, \pi)$ as the ``Mellin transform'' of the first Fourier coefficient of some appropriate $\phi(g)$ in the space of $\pi$. But the obvious problem is to choose this $\phi(g)$ or $W(g)$ correctly.

To gain some insight into this problem let us examine the special properties satisfied by $W_{\phi}(g)$ in case $\phi = \phi_f$. As in Section~\ref{ch:5}.B, it is easy to verify that the space of right translates of $W_{\phi}(g)$ generates a space of functions $W(\pi_f)$ on $G_{\A}$ satisfying the following properties:

\begin{enumerate}[(i)]

\item $W\!\left(\begin{bmatrix} 1 & x \\ 0 & 1 \end{bmatrix} g\right) = \tau(x)W(g)$ for all $x \in \A$;

\item $W\!\left(\begin{bmatrix} z & 0 \\ 0 & z \end{bmatrix} g\right) = \psi(z)W(g)$ for all $z \in \A^{\times}$;

\item $W$ is right K-finite, $C^{\infty}$ as a function of $G_{\infty}$, and rapidly decreasing; and

\item the natural representation of the Hecke algebra $\mathcal{H}(G_{\A})$ on $W(\pi_f)$ given by right convolution is equivalent to $\pi_f$.

\end{enumerate}

(Strictly speaking, $W(\pi_f)$ should be defined as the space of right convolutions of $W_{\phi}$ with elements of $\tilde{\mathcal{H}}(G_{\A})$ since the action of $G_{\A}$ at the infinite places does not preserve $K_{\infty}$-finiteness.)

This realization of the representation $\pi_f$ is called the Whittaker model for $\pi_f$ and $W(\pi_f)$ is called its Whittaker space. Thus $L(s,f)$ may be viewed described as the Mellin transform (along $\begin{bmatrix} y & 0 \\ 0 & 1 \end{bmatrix}$) of a certain privileged function in the Whittaker space of $\pi_f$.

What is special about $W_{\phi_f}(g)$ (among other things) is that it is right $K_0$ invariant. Thus, given an arbitrary constituent $\pi$ of $R_0^{\psi}$, one presumably could compute the conductor $c(\pi)$ of $\pi$ and select a ``correct'' $K_0^{c(\pi)}$ invariant function $W(g)$ in the Whittaker space of $\pi$. Then one could introduce $L(\pi,s)$ as some kind of Mellin transform of $W(g)$. This is in fact what Jacquet-Langlands do except that their attack on the problem gets closer to the heart of the matter. Jacquet-Langlands characterize special functions $W_v(g_v)$ in the Whittaker space of each local component $\pi_v$ of $\pi$ and then piece together these local functions to get $W(g)$; these local functions in turn are characterized in terms of the local ``zeta functions'' described in Subsection C below.

\section{Jacquet-Langlands' Theory}\label{sec:6c}

Suppose $\pi$ is any irreducible unitary (admissible) representation of $G_{\A} = \GL(2, \A)$. The purpose of this subsection is to explain how Jacquet-Langlands attach to each such $\pi$ an L-function $L(s, \pi)$ and how they characterize those $\pi$ that occur in $R_0^{\psi}$ in terms of the functions $L(s, \pi)$.

Analogous to Hecke's theory for $\GL(1)$ the first step will be to exploit the factorization $\pi = \otimes \pi_v$ announced in Theorem~\ref{thm:4.30}. (Cf. \eqref{eq:6.5}.) It should not seem surprising therefore that a ``global'' L-function $L(\pi, s)$ is defined by first introducing ``local L-factors'' $L(\pi_v, s)$ for each $\pi_v$ and then piecing these factors together. (Cf. the definition \eqref{eq:6.8}.) As already suggested, these local factors are described by introducing Whittaker models for each $\pi_v$ and then computing the Mellin transform of a certain privileged function in this space.

To make matters precise, fix once and for all a non-trivial additive character $\tau$ of $F \backslash \A$ and suppose $\tau(x) = \Pi\tau_v(x_v)$ for $x = (x_v)$ in $\A$.

\begin{theorem}[\normalfont (Existence and Uniqueness of the Local Whittaker Model)]\label{thm:6.8}

Suppose $\pi_v$ is an irreducible admissible (infinite-dimensional) representation of $G_v$, $v$ any place of $F$. Then in the space of locally constant functions on $G_v$ such that

\begin{equation}\label{eq:6.25}
W\!\left(\begin{bmatrix} 1 & x \\ 0 & 1 \end{bmatrix} g\right) = \tau_v(x)W(g)\tag{6.25}
\end{equation}

$(x \in F_v, g \in G_v)$ there is a unique subspace $W(\pi_v)$ which is invariant for the right action of $G_v$ and equivalent (as a $G_v$-module) to $\pi_v$.

\end{theorem}

As in the global case, $W(\pi_v)$ is called the Whittaker space of $\pi_v$, and $(G_v, W(\pi_v))$ its Whittaker model.

\begin{remark}\label{rem:6.9}

In case $v$ is archimedean some obvious modifications and additions are required. In this case we should consider irreducible admissible representations of the group algebra $\mathcal{H}(G_v)$ and consider only solutions of \eqref{eq:6.25} which are $C^{\infty}$ and such that

\begin{equation}\label{eq:6.26}
W\!\left(\begin{bmatrix} t & 0 \\ 0 & 1 \end{bmatrix}\right) = O(|t|)^N \text{ as } |t| \to \infty.\tag{6.26}
\end{equation}

Such solutions actually define classical Whittaker functions (along $\begin{bmatrix} t & 0 \\ 0 & 1 \end{bmatrix}$) (Hence the terminology for arbitrary fields.) The correct statement of Theorem~\ref{thm:6.8} for archimedean places is that there is a unique space of such functions on which the natural action of $\mathcal{H}(G_v)$ is equivalent to $\pi_v$.

\end{remark}

\begin{remark}[\normalfont (Concerning the proof of Theorem~\ref{thm:6.8})]\label{rem:6.10}

Jacquet-Langlands' proof of Theorem~\ref{thm:6.8} is simplest to explain in the archimedean case. In this case every $\pi_v$ is a subrepresentation of some $\rho(\mu_1, \mu_2)$ on $\mathcal{H}(\mu_1, \mu_2)$ (in the notation of Section~\ref{ch:4}). Thus (as in the proof of Theorem~\ref{thm:4.18}) one can introduce the space of functions

\begin{equation}\label{eq:6.27}
W_{\phi}(g) = \int \phi\!\left[ w^{-1}\begin{pmatrix} 1 & x \\ 0 & 1 \end{pmatrix} g \right] \overline{\tau_v(x)} \dd x\tag{6.27}
\end{equation}

with $\phi$ in the subspace of $\mathcal{H}(\mu_1, \mu_2)$ belonging to $\pi_v$. This space of functions can be shown to provide a Whittaker model for $\pi_v$. To establish uniqueness one uses the differential equations on p. 4.5 to show that such a space of functions has a basis each of whose elements satisfy a certain second order differential equation (coming from the Casimir operator pushed down to the symmetric space $G/K$) with only one solution (of sufficiently small growth at $\infty$! Cf. \eqref{eq:6.26}).

In the non-archimedean case a variant of the method of proof just sketched is required.

\end{remark}

The significance of Theorem~\ref{thm:6.8} for the theory of automorphic forms has already been suggested and should now be apparent. Indeed the local uniqueness of Whittaker models implies global uniqueness. Hence the multiplicity one result of Section~\ref{ch:5} follows immediately. On the other hand local existence also makes possible the choice of a privileged function in the space of $\pi_v$ and hence ultimately the choice of a function in the space of $\pi$ whose Mellin transform will coincide with $L(s, \pi)$. To explain this in more detail we need first to introduce the local and global ``zeta functions'' for $\pi_v$ and $\pi$.

\begin{definition}\label{def:6.11}

Let $\pi_v$ denote an irreducible admissible representation of $G_v$ and $W(\pi_v)$ its Whittaker space. Suppose $\chi$ is a unitary character of $(F_v)^{\times}$, $g \in G_v$, $W \in W(\pi_v)$, and $s$ a complex number. Then the local zeta function attached to $(g, \chi, W)$ is defined by the formula

\begin{equation}\label{eq:6.28}
\zeta(g, \chi, W, s) = \int_{F_v^{\times}} W\!\left(\begin{bmatrix} a & 0 \\ 0 & 1 \end{bmatrix} g\right) \chi(a) |a|^{s - 1/2} \dd^{\times} a.\tag{6.28}
\end{equation}

\end{definition}

\textbf{CAUTION:} Although $\xi(a) = W\!\left(\begin{bmatrix} a & 0 \\ 0 & 1 \end{bmatrix}g\right)$ is a relatively well behaved function on $F^{\times}$ it is not necessarily locally constant on $F$. Therefore the function of $s$ defined by \eqref{eq:6.28} does not coincide with the local zeta-function \eqref{eq:6.14} introduced for $\GL(1)$!

\begin{theorem}\label{thm:6.12}

\begin{enumerate}[(i)]

\item The integral defining $\zeta(g, \chi, W, s)$ converges for $s$ with sufficiently large real part;

\item There exists an Euler factor $L(s, \chi \otimes \pi_v)$ with the property that

\[
\frac{\zeta(g, \chi, W, s)}{L(s, \chi \otimes \pi_v)}
\]

is an entire function of $s$ for every $g, \chi$, and $W$, and such that

\begin{equation}\label{eq:6.29}
\zeta(1, \chi, W^0, s) = L(s, \chi \otimes \pi_v)\tag{6.29}
\end{equation}

for an appropriate choice of $W^0 \in W(\pi_v)$;

\item The function $\zeta(g, \chi, W, s)$ possesses an analytic continuation to the whole $s$-plane and satisfies the functional equation

\begin{equation}\label{eq:6.30}
\frac{\zeta(g, \chi, W, s)}{L(s, \chi \otimes \pi_v)} \varepsilon(\pi_v, \chi, s) = \frac{\zeta(wg, \chi^{-1}\psi_v^{-1}, W, 1 - s)}{L(1 - s, \psi_v^{-1}\chi_v^{-1} \otimes \pi_v)}\tag{6.30}
\end{equation}

where $\varepsilon(\pi_v, \chi, s)$ is independent of $g$ and $W$, $\psi_v$ is the central character of $\pi_v$, and $w = \begin{pmatrix} 0 & 1 \\ -1 & 0 \end{pmatrix}$.

\end{enumerate}

\end{theorem}

\begin{remark}[\normalfont (Concerning the statement of Theorem~\ref{thm:6.12})]\label{rem:6.13}

\begin{enumerate}[(i)]

\item By an Euler factor we understand a function of $s$ of the form $P^{-1}(q^s)$ where $P$ is a polynomial such that $P(0) = 1$ and $q = |\tilde{\omega}_v|^{-1}$ if $v$ is finite. For archimedean places $v$ the factor $L(s, \pi)$ will be a product of certain gamma functions to be specified below. In either case $L(s, \pi)$ is unique up to a constant factor. Indeed if $L(s, \pi)$ and $L^*(s, \pi)$ are two Euler factors satisfying the conditions of the theorem their quotient is an entire function without zeros.

\item The contragredient representation to $\pi_v$ is simply

\begin{equation}\label{eq:6.31}
\tilde{\pi}(g) = \psi_v^{-1}(g)\pi_v(g)\tag{6.31}
\end{equation}

if $\psi_v$ is the central character of $\pi_v$. Therefore \eqref{eq:6.30} reads

\begin{equation}\label{eq:6.32}
\frac{\zeta(W, s)}{L(s, \pi_v)} \varepsilon(\pi_v, s) = \frac{\zeta(\psi_v^{-1}W', 1 - s)}{L(1 - s, \tilde{\pi}_v)}\tag{6.32}
\end{equation}

if $g = 1$, $\chi \equiv 1$, and $W'$ denotes the right translate of $W$ by $\begin{pmatrix} 0 & 1 \\ -1 & 0 \end{pmatrix}$ (a kind of Fourier transform). The significance of \eqref{eq:6.32} is explained by recalling the local theory of Tate alluded to in Subsection A.

Tate's theory asserts that there is an Euler factor $L(s, \psi)$ such that (in the notation of \eqref{eq:6.14})

\[
\frac{\zeta(f_v, \psi_v, s)}{L(s, \psi)}
\]

is entire in $s$ for every $f \in S(F_v)$ and

\begin{equation}\label{eq:6.33}
\frac{\zeta(f_v, \psi_v, s)}{L(s, \psi_v)} \varepsilon(\psi_v, s) = \frac{\zeta(\hat{f}_v, \tilde{\psi}_v, 1 - s)}{L(1 - s, \tilde{\psi}_v)}\tag{6.33}
\end{equation}

for some (exponential) function $\varepsilon(\psi_v, s)$ which is independent of $s$. (Here $\overline{\psi}_v = \psi_v^{-1} =$ the ``contragredient'' of the representation $\psi_v$.) In fact some straightforward computation shows that

\begin{equation}\label{eq:6.34}
L(s, \psi_v) = (1 - \psi_v(\tilde{\omega})q^{-s})^{-1}\tag{6.34}
\end{equation}

if $\psi_v$ is unramified and

\begin{equation}\label{eq:6.35}
L(s, \psi_v) \equiv 1\tag{6.35}
\end{equation}

otherwise. In either case

\begin{equation}\label{eq:6.36}
L(s, \psi_v) = \zeta(f_v, \psi_v, s)\tag{6.36}
\end{equation}

for an appropriate choice of $f_v$ and the analogy between this situation and the situation for $\GL(2)$ described by Theorem~\ref{thm:6.12} is immediate.

For the record,

\[
\varepsilon(\psi, s) \equiv 1
\]

if $\psi$ is unramified and

\[
\varepsilon(\psi, s) = (\text{root of unity})\, q^{\operatorname{cond}(\psi)(\frac{1}{2} - s)}
\]

in general.

\item The factor $L(s, \chi \otimes \pi_v)$ may be interpreted as a g.c.d.\ for the family of meromorphic local zeta-functions $\zeta(g, W, \chi, s)$.

\item The special function $W_v^0(g)$ satisfying \eqref{eq:6.29} is the privileged local Whittaker function alluded to earlier. It will be the global function

\begin{equation}\label{eq:6.37}
W^0(g) = \prod_v W_v^0(g)\tag{6.37}
\end{equation}

whose adelic Mellin transform will coincide with $L(s,\pi)$.

\end{enumerate}

\end{remark}

\begin{theorem}\label{thm:6.14}

Suppose two irreducible admissible representations $\pi^1$ and $\pi^2$ of $G_v$ induce the same central character. Then $\pi^1$ and $\pi^2$ are equivalent if and only if

\begin{equation}\label{eq:6.38}
\varepsilon(\chi \otimes \pi^1, s) \frac{L(1 - s, \chi^{-1} \otimes \tilde{\pi}^1)}{L(s, \chi \otimes \pi^1)} = \varepsilon(\chi \otimes \pi^2, s) \frac{L(1 - s, \chi^{-1} \otimes \tilde{\pi}^2)}{L(s, \chi \otimes \pi^2)}\tag{6.38}
\end{equation}

for every character $\chi$ of $F_v^{\times}$.

\end{theorem}

This theorem says that the local factors $L(\chi \otimes \pi_v, s)$ and $\varepsilon(\chi \otimes \pi_v, s)$ uniquely determine $\pi_v$ among all the representations $\pi_v$ with given central character. Actually, given any $\pi_v$, these local factors can be explicitly described as follows.

\begin{theorem}[\normalfont (Description of $L(s,\pi)$ and $\varepsilon(s,\pi)$: The Non-Archimedean Case)]\label{thm:6.15}

In the table below we shall use the notation of Section~\ref{ch:4}; the factor $L(s, \psi)$ is the local factor for $\GL(1)$ given by \eqref{eq:6.34} and \eqref{eq:6.35}; finally

\[
E(\mu_1, \mu_2, s) = \frac{L(1 - s, \mu_1^{-1})}{L(s, \mu_2)} = \begin{cases} 1 & \text{if } \mu_1 \text{ is ramified} \\ -\mu_2(\tilde{\omega})q^{-s} & \text{otherwise} \end{cases}
\]

\textbf{TABLE}

{\small

\[
\begin{array}{@{}l@{\hspace{0.8em}}c@{\hspace{0.8em}}p{7.5cm}@{}}
\pi & L(s, \pi) & $\varepsilon(s, \pi)$\\ \hline
\pi = \pi(\mu_1, \mu_2) \ (\text{principal series}) & L(s, \mu_1)L(s, \mu_2) & $\varepsilon(s, \mu_1)\varepsilon(s, \mu_2)$ (equals 1 if $\mu_1$ and $\mu_2$ are unramified)\\
\pi = \sigma(\mu_1, \mu_2) \ (\text{special representation}) & L(\mu_1, s) & $\varepsilon(s, \mu_1)\varepsilon(s, \mu_2)E(\mu_1, \mu_2, s)$\\
\pi \ \text{supercuspidal} & 1 & unimportant to us
\end{array}
\]

}

\end{theorem}

\begin{theorem}[\normalfont (Description of $L(s,\pi)$ and $\varepsilon(s,\pi)$: The Archimedean Case)]\label{thm:6.16}

We restrict our attention to the case $F = \R$:

{\small

\[
\begin{array}{@{}p{5.4cm}@{\hspace{0.8em}}p{5.2cm}@{\hspace{0.8em}}p{5.6cm}@{}}
$\pi$ & $L(s, \pi)$ & $\varepsilon(s, \pi)$\\ \hline
$\pi = \pi(\mu_1, \mu_2)$ (principal series) & $L(s, \mu_1)L(s, \mu_2)$ & $\varepsilon(\mu_1, s)\varepsilon(\mu_2, s)$\\
$\pi = \sigma(p, t)$ (discrete series representation in $\pi(\mu_1, \mu_2)$) & $(2\pi)^{-s-\frac{t+p}{2}}\Gamma\!\left(s+\frac{t+p}{2}\right)$ & $i^{(p+1-n_1-n_2)}(-1)^{n_1+n_2}$\\[0.2em]
& & $(\mu_j(x) = |x|^{r_j}\mathrm{sgn}(x)^{n_j})$
\end{array}
\]

}

\end{theorem}

Here

\begin{equation}\label{eq:6.39}
L(s, \mu_j) = \pi^{-1/2(s + r_j + m_j)} \Gamma\!\left(\frac{s + r_j + m_j}{2}\right)\tag{6.39}
\end{equation}

and

\begin{equation}\label{eq:6.40}
\varepsilon(\mu_j, s) = i\tag{6.40}
\end{equation}

if $\mu_j(x) = |x|^{r_j}\mathrm{sgn}(x)^{m_j}$, $m_j = 0, 1$.

These tables will be helpful when we compare Jacquet-Langlands' theory with Hecke's classical theory. It will also be helpful to have the following information concerning the special function $W^0(g)$.

\begin{proposition}[\normalfont (Description of $W^{0}(g)$)]\label{prop:6.17}

\begin{enumerate}[(a)]

\item \emph{Archimedean case} $F_v = \R$: If $\pi_v$ belongs to the principal series for $G_v$ let $W_v^0(g)$ denote the unique $O(2)$-invariant function in $W(\pi_v)$; then $W_v^0(g)$ satisfies the relation

\begin{equation}\label{eq:6.41}
\zeta(1, \chi, W^0, s) = L(s, \chi \otimes \pi_v).\tag{6.41}
\end{equation}

If $\pi_v$ is equivalent to the discrete series representation $\sigma(\mathfrak{p}, \mathfrak{t})$ (in the notation of Remark~\ref{rem:4.6}) then \eqref{eq:6.41} is valid with $W_v^0(g)$ equal to the unique ``lowest weight vector'' in $W(\pi_v)$, i.e.\ the function in $W(\pi_v)$ such that $\rho(U)\phi = i(p + 1)\phi$ (cf.\ the notation of p. 57).

\item \emph{Non-archimedean} $F_v$: If $\pi_v$ is class one then $W(\pi_v)$ contains exactly one $K_v$ invariant function $W_v^0(g)$ such that $W_v^0(e) = 1$ and for this $W_v^0(g)$ \eqref{eq:6.29} obtains.

If $\pi_v$ is not class one then \eqref{eq:6.29} obtains with $W_v^0(g)$ equal to the function in $W(\pi_v)$ which is (almost) right-invariant for $K_p^{c(\pi_v)}$. See \eqref{eq:4.18} and \eqref{eq:4.19}.

\end{enumerate}

\end{proposition}

\textbf{CAUTION:} Parts of the last proposition are simply implicit in Chapter I of \cite{JacquetLanglands}.

We shall now explain how these local results are pieced together to obtain a global theory which generalizes the classical theory of Hecke.

Suppose $\pi$ is an irreducible unitary admissible representation of $\GL(2, \A)$. By Theorem~\ref{thm:4.30},

\[
\pi = \otimes \pi_v
\]

where almost every $\pi_v$ is class 1. On the other hand, every $\pi_v$ in \eqref{eq:6.42} possesses a unique Whittaker model $W(\pi_v)$ by Theorem~\ref{thm:6.8}, and the space of functions of $G_{\A}$ generated by functions of the form

\[
W(g) = \prod_v W_v(g_v),
\]

where $W_v = W_v^0$ for almost every $v$, provides a Whittaker model $W(\pi)$ for $\pi$. In particular the action of $\mathcal{H}(G_{\A})$ on $W(\pi)$ given by right convolution is equivalent to $\pi$. What is important is that this space $W(\pi)$ is unique. Indeed any space of functions on $G_{\A}$ satisfying properties (i)-(ii) of \eqref{eq:5.9} must by local uniqueness coincide with $W(\pi)$.

We remind the reader that if $\pi$ occurs in $R_0^{\psi}$ then an alternate construction of $W(\pi)$ results from the consideration of the first Fourier coefficients of functions in the space of $\pi$. More precisely, suppose $H^0$ denotes the space of K-finite functions in the space of $\pi$. The map

\[
\phi(g) \to \phi_1(g) = \int \phi\!\left(\begin{bmatrix} 1 & x \\ 0 & 1 \end{bmatrix} g\right) \overline{\tau(x)} \dd x
\]

then defines a one-to-one correspondence between $\pi$ and $W(\pi)$.

To state the main theorem of Jacquet-Langlands' theory we set

\begin{equation}\label{eq:6.42}
\varepsilon(\pi, \chi, s) = \prod_v \varepsilon(\pi_v \otimes \chi_v, s)\tag{6.42}
\end{equation}

for each grossencharacter $\chi = \Pi\chi_v$ of $F$. Since $\pi_v$ is almost always class 1 this product is actually a finite product by Theorem~\ref{thm:6.15}. We also put

\begin{equation}\label{eq:6.43}
L(s, \chi \otimes \pi) = \prod_v L(s, \chi_v \otimes \pi_v).\tag{6.43}
\end{equation}

This product converges for $\operatorname{Re}(s)$ sufficiently large since for almost every $v$, $L(s, \chi_v \otimes \pi_v)$ is of the form

\[
[1 - \mu_v\chi_v(\tilde{\omega}_v)q^{-s}]^{-1}[1 - \nu_v\chi_v(\tilde{\omega}_v)q^{-s}]^{-1}
\]

where $\mu_v$ and $\nu_v$ are local quasi-characters and $|\mu_v(x)| = |\nu_v(x)| = |x|^{\pm \sigma/2}$, $0 \leq \sigma \leq 1$.

\begin{theorem}\label{thm:6.18}

Suppose $\pi = \otimes \pi_v$ has central character $\psi$. Then $\pi$ occurs in $R_0^{\psi}$ if and only if $L(s, \chi \otimes \pi)$ satisfies the following properties for every grossencharacter $\chi$ of $F$:

\begin{enumerate}[(i)]

\item $L(s, \chi \otimes \pi)$ extends to an entire function bounded in vertical strips; and

\item $L(s, \chi \otimes \pi)$ satisfies the functional equation

\begin{equation}\label{eq:6.44}
L(s, \chi \otimes \pi) = \varepsilon(\pi, \chi, s) L(1 - s, \chi^{-1} \otimes \tilde{\pi})\tag{6.44}
\end{equation}

where $\tilde{\pi}(g) = \psi^{-1}(g)\pi(g)$.

\end{enumerate}

\end{theorem}

\begin{proof}[Proof (Sketch)]

Suppose first that $\pi$ occurs in $R_0^{\psi}$ and that its space is $H$. For each $v$ choose a function $W_v(g_v)$ in $W(\pi_v)$ such that $W_v = W_v^0$ almost everywhere. Then

\[
W(g) = \prod W_v(g_v)
\]

certainly belongs to $W(\pi)$ and

\begin{equation}\label{eq:6.45}
\phi(g) = \sum_{\xi \in F^{\times}} W\!\left(\begin{bmatrix} \xi & 0 \\ 0 & 1 \end{bmatrix} g\right)\tag{6.45}
\end{equation}

belongs to the space of K-finite functions in $H$.

Next consider the Mellin transform

\begin{equation}\label{eq:6.46}
\zeta(g, W, \chi, s) = \int_{F^{\times} \backslash \A^{\times}} \phi\!\left(\begin{bmatrix} x & 0 \\ 0 & 1 \end{bmatrix} g\right) \chi(x) |x|^{s - 1/2} \dd^{\times} x.\tag{6.46}
\end{equation}

Since $\phi\!\left(\begin{bmatrix} x & 0 \\ 0 & 1 \end{bmatrix} g\right)$ is rapidly decreasing at $\infty$ (and at $0$ by its right $G_F$-invariance) the integral in \eqref{eq:6.46} converges for all values of $s$ and represents an entire function bounded in vertical strips. But $W\!\left(\begin{bmatrix} x & 0 \\ 0 & 1 \end{bmatrix} g\right)$ also is rapidly decreasing at $\infty$. Therefore

\[
\int_{\A^{\times}} W\!\left(\begin{bmatrix} x & 0 \\ 0 & 1 \end{bmatrix} g\right) \chi(x) |x|^{s - 1/2} \dd^{\times} x = \prod_v \zeta(g_v, \chi_v, W_v, s)
\]

converges for $\operatorname{Re}(s)$ sufficiently large and in this range equals $\zeta(g, W, \chi, s)$. (Here $\zeta(g_v, \chi_v, W_v, s)$ is the local zeta function defined by \eqref{eq:6.28}.)

Now suppose $W_v(g)$ chosen so that

\[
\zeta(g_v, \chi_v, W_v, s) = L(s, \chi_v \otimes \pi_v).
\]

With this choice of $W(g) = \Pi W_v(g_v)$ we have (from above) that for $R(s)$ sufficiently large,

\[
\zeta(g, W, \chi, s) = \prod_v L(s, \chi_v \otimes \pi_v) = L(s, \chi \otimes \pi).
\]

Consequently $L(s, \chi \otimes \pi)$ obviously satisfies Condition (i) of Theorem 6.28 since $\zeta(g, \chi, W, s)$ does. To verify (ii) observe that

\[
\zeta(g, \chi, W, s) = \zeta(wg, \psi^{-1}\chi^{-1}, W, 1 - s)
\]

since $\phi(wg) = \phi(g)$. Consequently by the local functional equations,

\[
\begin{aligned}
\frac{\zeta(wg, \psi^{-1}\chi^{-1}, W, 1 - s)}{L(s, \chi^{-1} \otimes \tilde{\pi})} &= \prod_v \frac{\zeta(wg_v, \chi_v^{-1}\psi_v^{-1}, W_v, 1 - s)}{L(s, \chi_v^{-1} \otimes \tilde{\pi}_v)}\\
&= \prod_v \varepsilon(s, \chi_v, \pi_v) \frac{\zeta(g, \chi_v, W_v, s)}{L(s, \chi_v \otimes \pi_v)}\\
&= \varepsilon(s, \chi, \pi) \frac{\zeta(g, \chi, W, s)}{L(s, \chi \otimes \pi)}\\
&= \varepsilon(s, \chi, \pi) \frac{\zeta(wg, \psi^{-1}\chi^{-1}, W, 1 - s)}{L(s, \chi \otimes \pi)}
\end{aligned}
\]

and from this equality the functional equation \eqref{eq:6.44} is immediate.

Conversely, suppose $L(s, \chi \otimes \pi)$ satisfies Conditions (i) and (ii) of Theorem 6.28 for all $\chi$. To prove that $\pi$ then occurs in $R_0^{\psi}$ we simply ``reverse'' the procedure just outlined.

For each $W \in W(\pi)$ put

\begin{equation}\label{eq:6.47}
\phi_W(g) = \sum_{\xi \in F^{\times}} W\!\left(\begin{bmatrix} \xi & 0 \\ 0 & 1 \end{bmatrix} g\right).\tag{6.47}
\end{equation}

What has to be shown is that the map

\[
W \to \phi_W
\]

imbeds $W(\pi)$ into $L_0^2(G_F \backslash G_{\A}, \psi)$ and this reduces ultimately to showing that

\begin{equation}\label{eq:6.48}
\phi_W(wg) = \phi_W(g).\tag{6.48}
\end{equation}

Indeed the invariance of $\phi_W$ on the left by rational matrices of the form $\begin{bmatrix} a & b \\ 0 & c \end{bmatrix}$ is almost trivial from \eqref{eq:6.47} (assuming the convergence of \eqref{eq:6.47} has been established). Therefore Bruhat's decomposition together with known properties of $W(g)$ implies that \eqref{eq:6.48} is equivalent to $\phi$ belonging to $L_0^2(G_F \backslash G_{\A}, \psi)$.

To prove \eqref{eq:6.48} it certainly suffices to show that

\begin{equation}\label{eq:6.49}
\phi_W\!\left(\begin{bmatrix} x & 0 \\ 0 & 1 \end{bmatrix} g\right) = \phi_W\!\left(w\begin{bmatrix} x & 0 \\ 0 & 1 \end{bmatrix} g\right)\tag{6.49}
\end{equation}

for all $x \in \A^{\times}$. To prove \eqref{eq:6.49} it suffices in turn to verify that the Mellin transforms of the left and right hand sides are equal. Now the Mellin transform of the left hand side of \eqref{eq:6.49} is simply

\[
\begin{aligned}
\zeta(g, \chi, W, s) &= \int_{\A^{\times}} W\!\left(\begin{bmatrix} x & 0 \\ 0 & 1 \end{bmatrix} g\right) \chi(x) |x|^{s - \frac{1}{2}} \dd^{\times} x\\
&= \prod_v \zeta(g_v, \chi_v, W_v, s) \quad (\text{for } \operatorname{Re}(s) \gg 0)\\
&= L(s, \chi \otimes \pi) \prod_v \frac{\zeta(g_v, \chi_v, W_v, s)}{L(s, \chi_v \otimes \pi_v)},
\end{aligned}
\]

an entire function by virtue of the local theory and our assumptions on $L(s, \chi \otimes \pi)$. On the other hand the Mellin transform of the right hand side is

\[
\begin{aligned}
&\int_{F^{\times} \backslash \A^{\times}} \phi_W\!\left(w\begin{bmatrix} x & 0 \\ 0 & 1 \end{bmatrix} g\right) \chi(x) |x|^{s - \frac{1}{2}} \dd^{\times} x\\
&= \int \phi_W\!\left(\begin{bmatrix} x^{-1} & 0 \\ 0 & 1 \end{bmatrix} wg\right) \psi(x)\chi(x) |x|^{s - \frac{1}{2}} \dd^{\times} x\\
&= \int \phi_W\!\left(\begin{bmatrix} x & 0 \\ 0 & 1 \end{bmatrix} wg\right) \psi^{-1}(x)\chi^{-1}(x) |x|^{\frac{1}{2} - s} \dd^{\times} x\\
&= \zeta(wg, \psi^{-1}\chi^{-1}, W, 1 - s)\\
&= L(s, \chi^{-1} \otimes \tilde{\pi}) \prod_v \frac{\zeta(wg_v, \chi_v^{-1}\psi_v^{-1}, W_v, 1 - s)}{L(s, \chi_v^{-1} \otimes \tilde{\pi}_v)}.
\end{aligned}
\]

Therefore both Mellin transforms are entire functions and the fact that they coincide is an immediate consequence of the local and global functional equations.

\end{proof}

\begin{example}[\normalfont (Holomorphic Cusp Forms)]\label{exa:6.19}

Suppose $\pi_f = \otimes \pi_p$ is the representation of $\GL(2,\A)$ generated by the cusp form $f(z)$ in $S_k(\SL(2,\Z))$ whose eigenvalues for $T(p)$ are

\begin{equation}\label{eq:6.50}
a_p = p^{\frac{k-1}{2}}\left(p^{s_1} + p^{-s_1}\right).\tag{6.50}
\end{equation}

By Jacquet-Langlands' theory,

\[
\begin{aligned}
L(\pi, s) &= \prod_{p \leq \infty} L(s, \pi_p)\\
&= (2\pi)^{-s - \left(\frac{k-1}{2}\right)} \Gamma\!\left(s + \frac{k-1}{2}\right) \prod_{p < \infty} \left(1 - p^{-s_1}p^{-s}\right)^{-1}\left(1 - p^{s_1}p^{-s}\right)^{-1}
\end{aligned}
\]

(Cf. \eqref{eq:6.34}, Theorem  (\ref{thm:6.15}), and Theorem  (\ref{thm:6.16}).) Therefore by \eqref{eq:6.50},

\[
L(s, \pi) = (2\pi)^{-s - \left(\frac{k-1}{2}\right)} \Gamma\!\left(s + \frac{k-1}{2}\right) \prod \left(1 - p^{-s - \frac{k-1}{2}}a_p + p^{-2s}\right)^{-1}
\]

an L-function which agrees with Hecke's L-function

\[
L(s', f) = (2\pi)^{-s'}\Gamma(s')
\]

with $s' = s + \frac{k-1}{2}$. (Cf. \eqref{eq:1.28}.) Since in this case $\varepsilon(s, \pi) = i^k$, the functional equation \eqref{eq:6.44} is nothing but the functional equation \eqref{eq:1.27} of Hecke. (The substitution $s \to 1-s$ in \eqref{eq:6.44} corresponds to the substitution $s' \to k-s'$.) Consequently Jacquet-Langlands' theory is consistent with (at least this part of) Hecke's theory.

\end{example}

\begin{example}[\normalfont (Maass' Wave-forms)]\label{exa:6.20}

In \cite{Maass} certain L-functions are attached to real analytic cusp-forms in $W_s(\SL(2,\Z))$. As in Example~\ref{exa:6.19} it can be shown that these L-functions (and their functional equations) appear as special cases of the general theory sketched above.

\end{example}

\textbf{Concluding Remarks.} The ``real'' reason why the L-functions $L(s, \pi_f)$ and $L(s, f)$ of Example~\ref{exa:6.19} coincide is the following. For an appropriately chosen $W(g) = \Pi W_v(g)$ in $W(\pi_f)$, namely $W(g) = W_{\phi}(g)$ with $\phi = \phi_f$,

\begin{equation}\label{eq:6.51}
\begin{aligned}
L(s, \pi) &= \zeta(1, W, 1, s)\\
&= \int_{\A^{\times}} W\!\left(\begin{bmatrix} x & 0 \\ 0 & 1 \end{bmatrix}\right) |x|^{s - \frac{1}{2}} \dd^{\times} x.
\end{aligned}\tag{6.51}
\end{equation}

But since $W_v^0(g)$ is right $K_0$ invariant and transforms under $SO(2)$ according to the character $e^{-ik\theta}$,

\[
L(s, \pi) = \int_{\Q^{\times} \backslash \A^{\times}} \phi\!\left(\begin{bmatrix} y & 0 \\ 0 & 1 \end{bmatrix}\right) |y|^{s - \frac{1}{2}} \dd^{\times} y
\]

with

\[
\phi(g) = \phi_f(g) = f(g_{\infty}(i)) j(g_{\infty}, i)^{-k},
\]

or

\[
L(s, \pi) = \int_0^{\infty} f(iy) y^{s + \left(\frac{k-1}{2}\right) - 1} \dd y = L(s', f)!
\]

The equality \eqref{eq:6.51} furthermore emphasizes the fact that in proving the functional equation for $L(s,\pi)$, Jacquet-Langlands directly follow Hecke. Indeed it is (a generalization of) the Mellin transform which (once again) translates the automorphy condition for $\phi$ into a functional equation for $L(s,\pi_{\phi})$. In establishing their ``Converse to Hecke Theory'' Jacquet-Langlands follow Weil.

\section{Connections with the Classical Theory}\label{sec:6d}

Here we shall use Jacquet-Langlands' theory to derive at least one ``new'' result in the classical theory of forms.

The first natural question to ask is whether Theorem~\ref{thm:6.18} actually generalizes or simply overlaps Hecke's theory. In particular, does $L(s, \pi_f)$ always agree with Hecke's $L(s', f)!$ All we know thus far is that if $f$ is of level 1 then these functions do coincide.

So suppose now that $f$ belongs to $S_k(\Gamma_0(N))$. Then even in the context of the classical theory there is more than one way to attach an L-function to $f$. On the one hand, we have

\[
L(s, f) = \int_0^{\infty} f(iy) y^{s-1} \dd y.
\]

On the other, there is no loss of generality in assuming that

\begin{equation}\label{eq:6.52}
T(p)f = a_p f\tag{6.52}
\end{equation}

for all $p$ relatively prime to $N$. Then there are two possibilities. Either $f$ is a new form, in which case \eqref{eq:6.52} holds for all $p$, or $f$ is an old form, in which case $f(z) = h(dz)$ with $h$ a new form in $S_k(\Gamma_0(M))$ and

\[
T(p)h = a_p^{*} h \quad \text{for all } p.
\]

If the first possibility obtains, set

\[
L^*(s, f) = (2\pi)^{-s}\Gamma(s) \prod \left(1 - a_p p^{-s} + p^{k-1-2s}\right)^{-1}.
\]

If the second does, set

\[
L^*(s, f) = (2\pi)^{-s}\Gamma'(s) \prod_{p < \infty} \left(1 - a_p^{*} p^{-s} + p^{k-1-2s}\right)^{-1}.
\]

Note that $a_p^{*} = a_p$ for all $(p, N) = 1$. The natural question to ask is whether $L^*(s, f)$ coincides (up to a constant multiple) with $L(s, f)$ and the obvious answer is ``if and only if $f$ is a new form.''

\begin{proposition}\label{prop:6.21}

$L(s, \pi_f)$ and $L(s', f)$ coincide (up to a constant multiple) if and only if $f$ is a new form in $S_k(\Gamma_0(N))$.

\end{proposition}

\begin{proof}[Proof (Sketch)]

Suppose $f$ is a new form. Then by Proposition~\ref{prop:6.17}, and the discussion of Subsection C,

\[
\begin{aligned}
L(s, \pi_f) &= \int_{\Q^{\times} \backslash \A^{\times}} \phi_f\!\left(\begin{bmatrix} y & 0 \\ 0 & 1 \end{bmatrix}\right) |y|^{s - \frac{1}{2}} \dd y\\
&= \int_0^{\infty} f(iy) y^{s + \left(\frac{k-1}{2}\right) - 1} \dd y = L(s', f).
\end{aligned}
\]

On the other hand, suppose $f$ is an old form belonging to the new form $h$. Then

\[
L(s, \pi_f) = L(s, \pi_h) = L(s', h) \neq L(s', f).
\]

\end{proof}

The simple application of Jacquet-Langlands' theory we have in mind is reminiscent of Theorem~\ref{thm:5.14} whose representation-theoretic proof was promised earlier.

\textbf{Theorem~\ref{thm:5.14}.} Suppose $\pi^1$ and $\pi^2$ are two irreducible unitary constituents of $R_0^{\psi}$. Suppose also that the local components $\pi_v^1$ and $\pi_v^2$ are equivalent for all but finitely many finite places. Then $\pi^1 = \pi^2$.

\begin{proof}[Proof (Sketch)]

Let $S$ denote the finite set of finite places where it is not assumed $\pi_v^1 = \pi_v^2$. To prove that $\pi_v^1 = \pi_v^2$ for $p$ in $S$ we shall use the fact that the factors $L(\chi_v \otimes \pi_v, s)$ and $\varepsilon(\chi_v \otimes \pi_v, s)$ completely determine $\pi_v$ (Theorem~\ref{thm:6.14}).

Because $\pi^1$ and $\pi^2$ occur in $R_0^{\psi}$ the L-functions $L(s, \chi \otimes \pi^1)$ and $L(s, \chi \otimes \pi^2)$ satisfy functional equations of the type specified by Jacquet-Langlands' Theorem~\ref{thm:6.18}. Dividing the functional equation for $\pi^1$ by the functional equation for $\pi^2$, and rearranging slightly, we have

\begin{equation}\label{eq:6.53}
\begin{aligned}
&\prod_{v \in S} \frac{L(1 - s, \chi_v^{-1} \otimes \tilde{\pi}_v^1)\varepsilon(s, \chi_v \otimes \pi_v^1)}{L(s, \chi_v \otimes \pi_v^1)}\\
&= \prod_{v \in S} \frac{L(1 - s, \chi_v^{-1} \otimes \tilde{\pi}_v^2)\varepsilon(s, \chi_v \otimes \pi_v^2)}{L(s, \chi_v \otimes \pi_v^2)}.
\end{aligned}\tag{6.53}
\end{equation}

To conclude that $\pi_v^1 = \pi_v^2$ for each $v$ in $S$ one has only to eliminate all terms on each side of \eqref{eq:6.53} except the one involving $v$ and then apply Theorem~\ref{thm:6.14}. To do this one appeals, in turn, to the following two facts from \cite{JacquetLanglands}: (1) For any place $w$ of $F$ and any irreducible admissible representation $\pi_w$,

\[
L(s, \chi_w \otimes \pi_w) = 1
\]

and

\[
\varepsilon(s, \chi_w \otimes \pi_w) = \varepsilon(s, \chi_w\psi_w)\varepsilon(s, \chi_w)
\]

for $\chi_w$ a quasicharacter of $F_w^{\times}$ with sufficiently large conductor; and (2) For any character $\delta$ of $F_w^{\times}$ we can find a grossencharacter of $F$ which restricts to $\delta_u$ on $F_u^{\times}$ and has arbitrarily large conductor at the other places of $S$.

\end{proof}

The ``new'' result we have in mind is:

\begin{proposition}\label{prop:6.22}

Suppose $\pi^1$ and $\pi^2$ are two irreducible unitary constituents of $R_0^{\psi}$. Suppose also that the local components $\pi_p^1$ and $\pi_p^2$ are equivalent for all finite places. Then $\pi_{\infty}^1$ is equivalent to $\pi_{\infty}^2$ and $\pi^1 = \pi^2$. (We are working over $\Q$!)

\end{proposition}

\begin{proof}

By Theorem~\ref{thm:6.18} the L-functions $L(s, \pi^i)$ corresponding to $\pi^i$ satisfy the functional equation

\[
\frac{L(1 - s, \chi^{-1} \otimes \pi^i)}{L(s, \chi \otimes \pi^i)} \varepsilon(s, \chi \otimes \pi^i) = 1
\]

for all grossencharacters $\chi = \Pi\chi_p$ of $\Q$.

But by assumption, the quotients

\[
\frac{L(1 - s, \chi_p^{-1} \otimes \pi_p^i)\varepsilon(s, \chi_p \otimes \pi_p^i)}{L(s, \chi_p \otimes \pi_p)} \quad i = 1, 2
\]

coincide for all finite $p$. Therefore they coincide for $p = \infty$ as well, so by Theorem~\ref{thm:6.14}, $\pi_{\infty}^1 = \pi_{\infty}^2$.

\end{proof}

\begin{corollary}\label{cor:6.23}

Suppose $f(z)$ is a real analytic wave form for $\Gamma_0'(N)$ with $N$ square free. Then if

\[
T(p)f = a_p f \quad \text{for all } p,
\]

knowledge of $\{a_p\}$ completely determines the eigenvalue of $f$ for $\Delta$.

\end{corollary}

The special case of this corollary for $N = 1$ was communicated to me by Cartier before I had stated or proved Proposition~\ref{prop:6.22}. Its proof in general is similar to the proof of the following refinement of Hecke's Corollary~\ref{cor:1.10}.

\begin{theorem}\label{thm:6.24}

The Hecke ring has multiplicity one in the graded ring of holomorphic cusp forms of integral weight. More precisely, if

\[
f_i \in S_{k_i}(N, \psi) \text{ and } \quad T(p)f_i = a_p f_i \quad \text{for all } p,
\]

then $k_1 = k_2$, and $f_1$ is a multiple of $f_2$.

\end{theorem}

\begin{proof}

To give a representation-theoretic proof (i.e.\ to apply Proposition~\ref{prop:6.22}) we need to assume $N$ square free. In this case the representation $\pi_f$ corresponding to $f$ is determined at every finite place and hence at infinity as well. To prove the Theorem in general one argues directly with the functional equations for the Euler products

\[
(2\pi)^{-s}\Gamma(s) \prod_p \left(1 - a_p p^{-s} + p^{k-1-2s}\right)^{-1}
\]

after making the change of variables $k_i - 1$

\end{proof}

\section*{Notes and References}

\addcontentsline{toc}{subsection}{Notes and References}

\subsection*{Further Notes}

The analogy between Jacquet-Langlands' theory and Tate's proof of Theorem~\ref{thm:6.2} should not be overlooked. In fact it is possible to give another proof of the ``only if'' part of Theorem~\ref{thm:6.18} which completely avoids Whittaker models and follows Tate's proof directly rather than by analogy. This alternate approach is of considerable interest and has led to a fair amount of recent research which we shall at least sketch in the paragraphs below.

The basic idea is to introduce the local zeta functions

\[
\zeta(f_v, \pi_v, u, \tilde{u}, s) = \int_{G_v} f_v(x) \langle \pi_v(x)u, \tilde{u} \rangle |\det(x)|^{s + 1/2} \dd x
\]

for each Schwartz-Bruhat function $f_v$ on $M(2, F_v)$, each irreducible admissible representation $(\pi_v, V_v)$ of $G_v$, $u \in V_v$, $\tilde{u} \in \tilde{V}_v$, and $s \in \C$. These integrals can be shown to converge for $\operatorname{Re}(s)$ sufficiently large, to continue meromorphically to the whole $s$-plane, and to satisfy the functional equation

\begin{equation}
\frac{\zeta(f_v, \pi_v, u_v, \tilde{u}_v, s)}{L(s, \pi_v)} \varepsilon(s, \pi) = \frac{\zeta(\tilde{f}_v, \tilde{\pi}_v, u, \tilde{u}_v, 1 - s)}{L(1 - s, \tilde{\pi}_v)}.\tag{*}
\end{equation}

Here $\hat{f}_v$ denotes the (additive) Fourier transform of $f_v$, and $L(s,\pi)$ and $\varepsilon(s,\pi)$ are as in Theorem~\ref{thm:6.18}. The quotients

\[
\frac{\zeta(f_v, \pi_v, u_v, \tilde{u}_v, s)}{L(s, \pi_v)} = \Xi(s, f_v, u_v, \tilde{u}_v)
\]

are entire for all $f_v$ and

\[
\sum_{i=1}^{n} \Xi(s, f_v, u_i, \tilde{u}_i)
\]

is actually a non-zero constant for some choice of $f_v$, $u_1, \cdots, u_n$, and $\tilde{u}_1, \cdots, \tilde{u}_n$.

The above results extend Tate's local theory for $\GL(1)$ to $\GL(2)$ (cf. Section~\ref{ch:6}.A). What is remarkable is that the functions $\zeta(f_v, \pi_v, u, \tilde{u}, s)$ and $\zeta(g, 1, W, s)$ (cf. Definition~\ref{def:6.11}) lead to precisely the same Euler factors $L(s, \pi_v)$.

Now if $\pi = \bigotimes_v \pi_v$ is an irreducible admissible representation of $G_{\A}$ which occurs in $L_0^2(G_F \backslash G_{\A}, \psi)$ the global zeta-function

\[
\zeta(f, \pi, u, \tilde{u}, s) = \int_{G_{\A}} f(x) \langle \pi(x)u, \tilde{u} \rangle |\det(x)|^s \dd x
\]

is defined for $f \in S(M(2,\A))$, $u$ and $\tilde{u}$ in the dense subspace of $\psi$-cusp-forms, and $s$ in some half-plane. But by the Poisson summation formula, this function is shown to extend to an entire function in all of $\C$ which satisfies the functional equation

\[
\zeta(f, \pi, u, \tilde{u}, s) = \zeta(\tilde{f}, \tilde{\pi}, u, \tilde{u}, 2 - s)
\]

(cf. \eqref{eq:6.15}). Thus by proceeding just as Tate did for $\GL(1)$ one concludes (from the global and local theory) that the infinite product

\[
L(s, \pi) = \prod_v L(s, \pi_v)
\]

satisfies all the properties announced in Theorem~\ref{thm:6.18}.

This proof of the ``only if'' part of Theorem~\ref{thm:6.18} is to be found in Section 13 of \cite{JacquetLanglands}. Its generalization to arbitrary simple algebras is the subject matter of \cite{GodementJacquet}. The special case of division quaternion algebras is also described in Section 14 of \cite{JacquetLanglands} and Section~\ref{ch:10} of these Notes.

Yet another approach to the functional equation (*) is possible for certain operator-valued zeta functions attached to unitary representations. For the sake of definiteness, suppose that $G = \GL(2, \R)$ and $\pi$ is an irreducible unitary representation belonging to the discrete series for $G$.

If $f \in S(M(2, \R))$ the function $f(x)|\det(x)|^s$ belongs to $L^1(G)$ whenever $\operatorname{Re}(s) > 1$ and therefore the zeta-function integral

\[
\zeta(f, \pi, s) = \int_{G} f(x)\pi(x) |\det(x)|^s \dd x
\]

defines a holomorphic (operator-valued) function of $s$ in this range. The analytic continuation and functional equation of this integral is consistent with the above theory and results from computations in \cite{Gelbart} (in particular the relation (**) below) by a fairly direct argument. This argument, however, is not entirely obvious and has not been explained before. Therefore we shall sketch it below. We should also note that the complex analogue of the relation (**) below first appeared in \cite{Stein}.

Let $u$ and $v$ denote arbitrary vectors in the space of $\pi$. Then the integral

\[
\int_{G} f(g)\langle \pi(g)u, v \rangle |\det g|^s \dd g
\]

actually converges for $\operatorname{Re}(s) > \frac{1}{2}$. To see this it suffices to integrate over $G^{+} = \GL^{+}(2, \R) = \R_{+}^{\times} \times \SL(2, \R)$. But by Cauchy-Schwarz and the square-integrability of $\langle \pi(g)u, v \rangle$,

\[
\int_{G^{+}} f(g)\langle \pi(g)u, v \rangle |\det g|^s \dd g \leq M \left\{\int_{\SL(2, \R)} |f_s(x)|^2 \dd x \right\}^{1/2}
\]

where $f_s(x) = \int_0^{\infty} f(xr) r^s \frac{\dd r}{r}$, so to prove our claim it suffices to establish the following:

\textbf{Lemma.} Let $S = \SL(2, \R)$. Then $f_s \in L^2(S)$ when $\operatorname{Re}(s) > \frac{1}{2}$.

\begin{proof}[Proof (Sketch)]

Because

\[
\int_{G^{+}} f(g)|\det g|^s \dd g = \int_{S} f_s(x) \dd x,
\]

it follows that $\|f_S\|_1 < \infty$ when $\operatorname{Re}(s) = 1 + \varepsilon, \varepsilon > 0$. On the other hand, $\|f_S\|_{\infty} < \infty$ whenever $\operatorname{Re}(s) = \varepsilon > 0$. So combining these facts with the Hadamard 3-line theorem we conclude that for $\operatorname{Re}(s) > \frac{1}{2}$,

\[
\int_{S} f_s(x)\phi(x) \dd x \leq A\|\phi\|_2
\]

for all $\phi$ in a dense subspace of $L^2(S)$. The desired result then follows from the converse to H\"older's inequality.

\end{proof}

Now set

\[
\zeta(\hat{f}, \widetilde{\pi}, 2 - s) = \int \hat{f}(g)\langle \widetilde{\pi}(g)u, v \rangle |\det g|^{2 - s} \dd g
\]

where $\widetilde{\pi}(g)$ denotes the discrete series representation $\pi(g^{-1})$. By the above argument this integral converges for $\operatorname{Re}(s) < \frac{3}{2}$.

\textbf{Theorem.} The zeta-function $\zeta(f, \pi, s)$, initially defined for $\operatorname{Re}(s) > \frac{1}{2}$, continues meromorphically to the whole $s$-plane and satisfies the functional equation

\[
\zeta(\hat{f}, \widetilde{\pi}, 2 - s) = \mathsf{m}_{\pi}(s)\zeta(f, \pi, s)
\]

where

\[
\mathsf{m}_{\pi}(s) = i^{n+1}(2\pi)^{s + it - 1} \frac{\Gamma'\left(\frac{n + (2 - s) - it}{2}\right)}{\Gamma'\left(\frac{n + s + it}{2}\right)}
\]

if $\pi$ is indexed by the parameters $n$ and $it$. (Cf. Remark~\ref{rem:4.6}.)

\begin{proof}

The integrals defining $\zeta(f, \pi, s)$ and $\zeta(f, \tilde{\pi}, 2-s)$ converge simultaneously in the strip

\[
\frac{1}{2} < \operatorname{Re}(s) < \frac{3}{2}
\]

and so define analytic functions there. By Theorem 3 of \cite{Gelbart}

\begin{equation}
\int_{G} \hat{f}(x)\widetilde{\pi}(g) |g| \dd g = \mathsf{m}_{\pi}(1) \int_{G} f(g)\pi(g) |g| \dd g\tag{**}
\end{equation}

which says that

\[
\zeta(\hat{f}, \widetilde{\pi}, 2 - s) = \mathsf{m}_{\pi}(s)\zeta(f, \pi, s)
\]

when $\operatorname{Re}(s) = 1$. By analytic continuation then this relation also holds throughout the strip $\frac{1}{2} < \operatorname{Re}(s) < \frac{3}{2}$ and hence everywhere since $\mathsf{m}_{\pi}(s)$ is obviously meromorphic in $\C$.

\end{proof}

An alternate approach to this theorem due to Godement is described in Section~\ref{ch:9} or \cite{GodementJacquet}. See also Section 13 of \cite{JacquetLanglands}. The special case of the trivial representation is discussed in \cite{Weil6} as well as in the above references.

\subsection*{Further References}

\textbf{A.} A good reference for Hecke's theory of L-series with grossencharacter is \cite{Godement3}. Tate's thesis appears in \cite{CasselsFrohlichs} but is also discussed in [Lang] and \cite{Goldstein}.

\textbf{B.} An essential feature of Jacquet-Langlands' theory is that it puts on a completely equal footing the real analytic wave-forms of Maass and the classical holomorphic cusp forms of Hecke's theory. This feature as well as the generality of an arbitrary base field is shared by the elementary theory of \cite{Weil2} which runs parallel to Jacquet-Langlands theory without using the theory of group representations.

\textbf{C.} The results quoted here are (sometimes implicitly) proved in \cite{JacquetLanglands}.

The fundamental local theorem on the existence and uniqueness of Whittaker models (Theorem~\ref{thm:6.8}) was used by Jacquet-Langlands to classify all irreducible admissible representations as follows.

For each $W \in W(\pi_v)$ set

\[
f_W(x) = W\!\left(\begin{bmatrix} x & 0 \\ 0 & 1 \end{bmatrix}\right).
\]

The resulting space of functions on $F_v^{\times}$ is called the Kirillov space of $\pi_v$ and its properties are described by the following reformulation of the existence and uniqueness theorem for local Whittaker models.

\textbf{Theorem.} Every irreducible admissible representation $\pi_v$ of $G_v$ (of infinite dimension) is equivalent to a unique representation acting in a space of locally constant functions on $F_v^{\times}$ compactly supported in $F_v$ and such that

\begin{equation}\label{eq:6.54}
\pi_v\!\left(\begin{bmatrix} a & b \\ 0 & 1 \end{bmatrix}\right) f(x) = \tau_v(bx)f(ax) \quad \left(a \in F_v^{\times}, b \in F_v\right).\tag{6.54}
\end{equation}

This space $K(\pi_v)$ furthermore contains the Schwartz-Bruhat space $S(F_v^{\times})$ as a subspace of codimension at most two.

One passes from the Kirillov model to the Whittaker model by setting

\[
W_f(g) = (\pi_v(g)f)(1).
\]

The classification of the irreducible admissible (infinite dimensional) representation of $G_v$ then takes the form that $\pi_v$ is supercuspidal if $K(\pi_v) = S(F_v^{\times})$, is special if $S(F_v^{\times})$ has codimension 1, and belongs to the principal series otherwise (i.e.\ if $S(F_v^{\times})$ has codimension two in $K(\pi_v)$). Note that the degree of $L(s, \pi_v)$ is at most the codimension of $S(F_v^{\times})$ in $K(\pi_v)$.

We remark also that by virtue of \eqref{eq:6.54} an irreducible representation $\pi_v$ is completely determined by knowledge of the single operator $\pi\!\left(\begin{pmatrix} 0 & 1 \\ -1 & 0 \end{pmatrix}\right)$ once one knows the central character of $\pi_v$ and its space $K(\pi_v)$.

The terminology for $K(\pi_v)$ seems apt because it was Kirillov who first proved that (over any local field) every irreducible unitary representation $\pi_v$ of $G_v$ remains irreducible upon restriction to the subgroup

\[
\left\{\begin{bmatrix} a & b \\ 0 & 1 \end{bmatrix}\right\}
\]

and on this subgroup is of the form \eqref{eq:6.54}. In the real case such facts had already been exploited in \cite{KunzeStein}.

A very readable account of Theorem~\ref{thm:6.8} is to be found in \cite{Godement}. Using the theory of spherical functions a completely different argument for uniqueness has recently been described by R.~Howe.
